\documentclass[12pt,a4paper]{article}

\usepackage[utf8]{inputenc}
\usepackage[T5]{fontenc}
\usepackage[vietnamese]{babel}
\usepackage{geometry}
\usepackage{hyperref}
\usepackage{longtable}
\usepackage{booktabs}
\usepackage{array}
\usepackage{enumitem}
\usepackage{listings}
\usepackage{xcolor}
\usepackage{graphicx}

\geometry{margin=2.4cm}
\hypersetup{
  colorlinks=true,
  linkcolor=blue!60!black,
  urlcolor=blue!60!black,
  pdftitle={Bao cao du an GraphRAG},
  pdfauthor={GitHub Copilot},
}

\definecolor{codebg}{RGB}{248,248,248}
\definecolor{codeborder}{RGB}{220,220,220}

\lstdefinestyle{codeblock}{
  backgroundcolor=\color{codebg},
  basicstyle=\ttfamily\small,
  frame=single,
  rulecolor=\color{codeborder},
  breaklines=true,
  columns=fullflexible,
  showstringspaces=false
}

\title{\textbf{Báo cáo chi tiết dự án GraphRAG}\large Hệ thống hỏi đáp tài liệu khoa học dựa trên Knowledge Graph}
\author{GitHub Copilot}
\date{\today}

\begin{document}

\maketitle
\tableofcontents
\newpage

\section{Tổng quan dự án}
Dự án \textbf{GraphRAG} này là một hệ thống truy vấn và suy luận trên tập tài liệu khoa học, kết hợp giữa \textit{Retrieval-Augmented Generation} và \textit{Knowledge Graph}. Thay vì chỉ chia tài liệu thành các đoạn văn bản rồi tìm kiếm theo vector, hệ thống trích xuất cấu trúc của bài báo thành các thực thể và quan hệ có ngữ nghĩa như \textit{Paper}, \textit{Author}, \textit{Organization}, \textit{Conference}, \textit{Topic}, \textit{Methodology}, \textit{Dataset}, \textit{Result}.

Mục tiêu cốt lõi của dự án là:
\begin{itemize}[leftmargin=1.6em]
  \item Nạp PDF khoa học vào hệ thống một cách tự động.
  \item Trích xuất metadata và cấu trúc tri thức từ nội dung tài liệu.
  \item Lưu trữ tri thức vào Neo4j dưới dạng property graph.
  \item Hỗ trợ người dùng hỏi đáp bằng ngôn ngữ tự nhiên thông qua agent đa bước.
  \item Cung cấp giao diện web để tải lên tài liệu, xem graph, chat và duyệt tài liệu.
\end{itemize}

Ở mức kiến trúc, dự án gồm 5 khối lớn: \textbf{Frontend}, \textbf{Backend API}, \textbf{Worker xử lý nền}, \textbf{Knowledge Graph/Database}, và \textbf{Pipeline xử lý tài liệu}. Toàn bộ được đóng gói bằng Docker Compose để thuận tiện triển khai nội bộ và demo.

\section{Mục tiêu và bài toán}
Trong hệ thống RAG truyền thống, tài liệu thường được chia nhỏ thành chunk và truy hồi theo độ tương đồng ngữ nghĩa. Cách này nhanh nhưng có hạn chế:
\begin{itemize}[leftmargin=1.6em]
  \item Mất cấu trúc quan hệ giữa các thực thể trong bài báo.
  \item Khó truy vấn các câu hỏi có tính so sánh hoặc quan hệ đa bước.
  \item Khó thống nhất cùng một thực thể khi xuất hiện ở nhiều bài.
\end{itemize}

GraphRAG giải quyết bằng cách lưu tri thức theo đồ thị. Khi người dùng hỏi câu như ``paper nào dùng Transformer và đánh giá trên SQuAD?'' hệ thống không chỉ truy theo text mà còn có thể đi qua cạnh \textit{USES\_METHOD}, \textit{EVALUATED\_ON}, \textit{AUTHORED}, \textit{AFFILIATED\_WITH}, \textit{CITES} để suy luận có cấu trúc hơn.

\section{Tổng quan kiến trúc hệ thống}
Kiến trúc tổng thể được mô tả như sau:

\begin{lstlisting}[style=codeblock,caption={Luồng kiến trúc mức cao}]
User -> Next.js Frontend -> FastAPI Backend
     -> Celery Worker -> PDF Parsing -> LLM Extraction
     -> Neo4j Knowledge Graph + PostgreSQL Metadata
     -> LangGraph Agent -> Answer Synthesis
\end{lstlisting}

Hệ thống hiện tại chạy trên Docker Compose với các service chính:
\begin{itemize}[leftmargin=1.6em]
  \item \textbf{Neo4j}: lưu đồ thị tri thức.
  \item \textbf{PostgreSQL}: lưu metadata tài liệu, trạng thái upload, lịch sử chat.
  \item \textbf{Redis}: hàng đợi cho Celery.
  \item \textbf{Backend FastAPI}: API trung tâm.
  \item \textbf{Celery worker}: xử lý PDF, trích xuất entity/relation, ghi vào graph.
  \item \textbf{Frontend Next.js}: giao diện tải tài liệu, chat, và trực quan graph.
  \item \textbf{Nginx}: reverse proxy.
\end{itemize}

\subsection{Sơ đồ dữ liệu và luồng xử lý}
Luồng ingest tài liệu đi theo chuỗi:
\begin{enumerate}[leftmargin=1.6em]
  \item Người dùng upload PDF.
  \item Backend lưu file và tạo bản ghi metadata.
  \item Celery worker nhận tác vụ ingest.
  \item PDF được parse bằng PyMuPDF.
  \item Metadata được trích xuất từ phần đầu tài liệu bằng LLM.
  \item Tài liệu được chia section, sau đó đi qua trích xuất entity/relation.
  \item Các thực thể/quan hệ được ghi vào Neo4j.
  \item Embedding và FAISS được cập nhật để phục vụ tìm kiếm vector.
\end{enumerate}

Luồng chat đi theo:
\begin{enumerate}[leftmargin=1.6em]
  \item Người dùng nhập câu hỏi.
  \item Backend đưa câu hỏi vào LangGraph agent.
  \item Planner sinh kế hoạch truy vấn.
  \item Retriever sinh Cypher và query Neo4j.
  \item Synthesizer tổng hợp ngữ cảnh và trả lời cuối.
\end{enumerate}

\section{Thành phần chính của hệ thống}
\subsection{Frontend}
Frontend được xây dựng bằng Next.js và đóng vai trò giao diện duy nhất cho người dùng cuối. Các màn hình chính gồm:
\begin{itemize}[leftmargin=1.6em]
  \item Trang upload tài liệu.
  \item Trang papers để duyệt danh sách tài liệu.
  \item Trang chat để hỏi đáp với agent.
  \item Trang graph để trực quan hóa knowledge graph.
\end{itemize}

Điểm quan trọng là trang graph không chỉ hiển thị node/edge, mà còn cho phép lọc, chọn, hover và quan sát mối liên kết giữa các thực thể. Trong quá trình hoàn thiện gần đây, graph đã được chuyển sang renderer tương tác hơn để tránh tình trạng label chồng lên nhau và hỗ trợ kéo/zoom trực tiếp.

\subsection{Backend FastAPI}
Backend là trung tâm điều phối API. Vai trò chính gồm:
\begin{itemize}[leftmargin=1.6em]
  \item Nhận yêu cầu upload và chat.
  \item Quản lý metadata tài liệu.
  \item Cung cấp API graph từ Neo4j.
  \item Điều phối Celery worker cho các tác vụ nặng.
\end{itemize}

Backend còn là nơi đặt các lớp core như client Neo4j, client LLM, database, và các model Pydantic/SQLAlchemy.

\subsection{Celery worker}
Worker là nơi chạy các tác vụ tiêu tốn thời gian như parse PDF, gọi LLM, trích xuất entity/relation, ghi đồ thị, và cập nhật chỉ mục FAISS. Đây là phần quyết định chất lượng dữ liệu đầu vào của toàn hệ thống.

Trong thực tế sử dụng, worker cần được rebuild khi thay đổi code vì image của worker được build từ Dockerfile, không mount source trực tiếp.

\subsection{Neo4j knowledge graph}
Neo4j lưu các node và relationship theo mô hình property graph. Các node chính bao gồm paper, author, organization, conference, topic, methodology, dataset, result, task. Các edge mô tả quan hệ như authored, affiliated_with, published_at, uses_method, evaluates_on, cites, achieves.

Đây là tầng tri thức giúp hệ thống trả lời các câu hỏi cần tính liên kết giữa nhiều paper hoặc nhiều thực thể khác nhau.

\subsection{PostgreSQL metadata store}
PostgreSQL lưu thông tin vận hành của tài liệu và hội thoại, ví dụ:
\begin{itemize}[leftmargin=1.6em]
  \item Tên file.
  \item Trạng thái xử lý.
  \item Số lượng entity/relation đã trích xuất.
  \item Lịch sử chat và session.
\end{itemize}

\subsection{FAISS vector store}
Song song với knowledge graph, hệ thống vẫn duy trì index vector trong FAISS. Vai trò của FAISS là hỗ trợ tìm kiếm ngữ nghĩa theo chunk văn bản. Điều này bổ sung cho graph: graph phục vụ suy luận quan hệ, còn vector search phục vụ truy hồi ngôn ngữ tự nhiên và ngữ cảnh cục bộ.

\section{Pipeline xử lý tài liệu}
Pipeline tài liệu là phần cốt lõi của dự án. Nó biến PDF thành ba loại dữ liệu:
\begin{enumerate}[leftmargin=1.6em]
  \item Metadata của paper.
  \item Knowledge graph nodes/edges.
  \item Embeddings/chunks cho truy hồi vector.
\end{enumerate}

\subsection{Parse PDF}
Giai đoạn đầu dùng PyMuPDF để đọc PDF, lấy full text, số trang và các đoạn đầu tài liệu. Parser có xu hướng hỗ trợ layout hai cột, vì nhiều paper khoa học không nằm trên một cột đơn giản.

\subsection{Trích xuất metadata}
Metadata gồm title, authors, year, abstract, keywords, venue, doi, categories. Hiện tại hệ thống đã tránh dùng structured function-calling trực tiếp cho metadata vì một số model LLM có thể trả về lỗi tool-calling. Thay vào đó, response được ép về JSON thô rồi parse lại.

Đây là một điểm quan trọng vì metadata là đầu vào cho các bước sau, đặc biệt là nối author với paper và suy ra tổ chức của tác giả.

\subsection{Chia section và chunking}
Sau khi parse, tài liệu được chia thành các section. Mỗi section tiếp tục được chunk để phục vụ embedding. Cần nhấn mạnh rằng chunking trong dự án này phục vụ vector search; nó \textbf{không thay thế} entity/relation extraction.

\subsection{Embedding và cập nhật FAISS}
Các chunk được mã hóa bằng SentenceTransformers, rồi ghi vào FAISS. Việc này cho phép tìm kiếm theo ngữ nghĩa trong kho tài liệu nội bộ.

\subsection{Ghi vào Neo4j}
Sau khi có metadata và extraction, hệ thống ghi paper, author, organization, conference, topic, methodology, dataset, result và các relationship vào Neo4j.

Các cạnh quan trọng nhất trong luồng hiện tại là:
\begin{itemize}[leftmargin=1.6em]
  \item \texttt{AUTHORED}: Author $\rightarrow$ Paper.
  \item \texttt{AFFILIATED\_WITH}: Author $\rightarrow$ Organization.
  \item \texttt{PUBLISHED\_AT}: Paper $\rightarrow$ Conference.
  \item \texttt{USES\_METHOD}: Paper $\rightarrow$ Methodology.
  \item \texttt{EVALUATED\_ON}: Paper $\rightarrow$ Dataset.
  \item \texttt{ACHIEVES}: Paper $\rightarrow$ Result.
\end{itemize}

\section{Mô hình dữ liệu tri thức}
Graph schema được mô tả trong tài liệu dự án và tuân theo mô hình labeled property graph.

\subsection{Các loại node}
\begin{longtable}{>{\raggedright\arraybackslash}p{3.2cm} p{10.6cm}}
\toprule
\textbf{Node} & \textbf{Ý nghĩa} \\
\midrule
Paper & Đại diện cho một bài báo khoa học; node trung tâm của graph. \\
Author & Tác giả hoặc đồng tác giả của paper. \\
Organization & Trường, viện, phòng thí nghiệm, hoặc đơn vị công tác. \\
Conference & Hội nghị, tạp chí, hoặc venue xuất bản. \\
Topic & Chủ đề nghiên cứu ở mức khái quát. \\
Task & Bài toán cụ thể mà paper giải quyết. \\
Methodology & Phương pháp, mô hình, hoặc kỹ thuật được sử dụng. \\
Dataset & Tập dữ liệu dùng để huấn luyện/đánh giá. \\
Result & Kết quả metric cụ thể đạt được. \\
Chunk & Đoạn văn bản bổ trợ cho vector search, không thay thế graph extraction. \\
\bottomrule
\end{longtable}

\subsection{Các loại relationship}
\begin{longtable}{>{\raggedright\arraybackslash}p{4cm} p{8cm}}
\toprule
\textbf{Quan hệ} & \textbf{Ý nghĩa} \\
\midrule
AUTHORED & Author viết Paper. \\
AFFILIATED\_WITH & Author thuộc Organization. \\
PUBLISHED\_AT & Paper được xuất bản tại Conference. \\
COVERS\_TOPIC & Paper thuộc về một Topic. \\
USES\_METHOD & Paper sử dụng một Methodology. \\
EVALUATED\_ON & Paper được đánh giá trên Dataset. \\
CITES & Paper trích dẫn một Paper khác. \\
ACHIEVES & Paper đạt được một Result. \\
SUBTOPIC\_OF & Topic con thuộc Topic cha. \\
VARIANT\_OF & Methodology là biến thể của Methodology khác. \\
\bottomrule
\end{longtable}

\subsection{Vì sao graph quan trọng}
Graph giúp trả lời các câu hỏi như:
\begin{itemize}[leftmargin=1.6em]
  \item Paper nào cùng dùng một phương pháp?
  \item Những tác giả nào cùng thuộc một tổ chức?
  \item Paper nào công bố ở cùng venue?
  \item Paper nào vừa dùng một dataset cụ thể vừa đạt một metric cụ thể?
\end{itemize}

\section{Agent và luồng suy luận}
Hệ thống sử dụng LangGraph để điều phối agent. Mô hình agent hiện tại gồm ba vai trò chính:
\begin{enumerate}[leftmargin=1.6em]
  \item \textbf{Planner}: phân tích câu hỏi và tạo kế hoạch nhiều bước.
  \item \textbf{Retriever}: sinh truy vấn Cypher và lấy dữ liệu từ Neo4j.
  \item \textbf{Synthesizer}: tổng hợp kết quả trả lời cho người dùng.
\end{enumerate}

Ưu điểm của LangGraph là cho phép agent có nhiều vòng lặp và có thể mở rộng thêm validation hoặc fallback sau này.

\subsection{Kịch bản truy vấn mẫu}
Ví dụ người dùng hỏi:
\begin{quote}
  ``Những paper nào dùng Transformer và đánh giá trên SQuAD?''
\end{quote}

Luồng xử lý sẽ là:
\begin{itemize}[leftmargin=1.6em]
  \item Phân tích câu hỏi để xác định entity cần tìm.
  \item Query Neo4j theo \texttt{USES\_METHOD} và \texttt{EVALUATED\_ON}.
  \item Gom context và trả lời bằng ngôn ngữ tự nhiên.
\end{itemize}

\section{Giao diện người dùng}
Frontend hiện tại là một phần rất quan trọng của dự án vì đây là nơi người dùng cảm nhận được tri thức từ đồ thị.

\subsection{Trang upload}
Cho phép người dùng tải PDF lên hệ thống. Sau khi upload, file được đẩy vào pipeline và worker bắt đầu xử lý nền.

\subsection{Trang papers}
Hiển thị danh sách tài liệu đã nhập, cho phép xem nhanh metadata và trạng thái xử lý.

\subsection{Trang chat}
Đây là giao diện truy vấn tri thức. Chat page kết nối tới backend để trả lời dựa trên data trong graph.

\subsection{Trang graph}
Trang graph là màn hình trực quan hóa Knowledge Graph. Nó cho phép:
\begin{itemize}[leftmargin=1.6em]
  \item Lọc node theo từ khóa.
  \item Chọn node để xem chi tiết.
  \item Hover node để xem label.
  \item Kéo node và zoom graph.
\end{itemize}

Trong quá trình hoàn thiện gần đây, graph UI đã được sửa từ cách vẽ SVG thủ công sang renderer tương tác hơn để giải quyết vấn đề label chồng lấn và graph quá rối khi số node tăng.

\section{Các vấn đề thực tế đã gặp và hướng xử lý}
Trong quá trình chạy dự án, một số vấn đề kỹ thuật quan trọng đã xuất hiện.

\subsection{Lỗi metadata trả về \texttt{Unknown}}
Khi dùng structured-output trực tiếp với Groq, API có thể trả lỗi tool-calling. Hệ quả là metadata fallback về \texttt{Unknown}. Cách xử lý là chuyển sang đọc raw JSON text và tự parse thay vì ép structured tool-call ở metadata extraction.

\subsection{Chunking bị hiểu nhầm là logic chính}
Trong đồ án, chunking chỉ phục vụ vector search và embedding. Nó không phải logic chính của graph extraction. Vì vậy cần phân biệt rõ:
\begin{itemize}[leftmargin=1.6em]
  \item Chunking: phục vụ FAISS/vector retrieval.
  \item Entity/relation extraction: phục vụ Neo4j knowledge graph.
\end{itemize}

\subsection{Node tổ chức bị đứng riêng}
Một số node như \textit{Google Brain} ban đầu xuất hiện như node mồ côi. Nguyên nhân không phải frontend, mà là pipeline chưa tạo đủ cạnh \texttt{AFFILIATED\_WITH}. Để xử lý, hệ thống đã bổ sung heuristic suy luận tổ chức từ header PDF và backfill cạnh author--organization vào Neo4j.

\subsection{Worker cần rebuild sau thay đổi code}
Do Celery worker được build thành image riêng, thay đổi code trong pipeline không có hiệu lực nếu chỉ restart container. Cần dùng:
\begin{lstlisting}[style=codeblock]
docker compose up -d --build backend celery-worker
\end{lstlisting}

\section{Đánh giá hệ thống}
\subsection{Điểm mạnh}
\begin{itemize}[leftmargin=1.6em]
  \item Kiến trúc rõ ràng, tách biệt ingest, graph, và chat.
  \item Có cả graph retrieval và vector retrieval.
  \item Hỗ trợ ingest PDF khoa học tự động.
  \item Đã có nền tảng cho multi-hop reasoning qua LangGraph.
  \item UI graph cho người dùng nhìn được cấu trúc tri thức.
\end{itemize}

\subsection{Hạn chế hiện tại}
\begin{itemize}[leftmargin=1.6em]
  \item Một số module backend vẫn còn là skeleton hoặc chưa hoàn thiện đầy đủ.
  \item Extraction phụ thuộc mạnh vào chất lượng LLM.
  \item Một số PDF layout phức tạp vẫn có thể làm lệch extraction.
  \item Graph schema và query layer cần tinh chỉnh thêm cho đồng nhất.
\end{itemize}

\subsection{Độ sẵn sàng demo}
Xét ở góc độ demo, hệ thống đã có các thành phần chính: upload, ingest, graph visualization, chat pipeline, và knowledge graph lưu trữ. Tuy nhiên để đưa vào vận hành ổn định hơn, nên tiếp tục cải thiện fallback heuristic, backfill edge, và khả năng xử lý lỗi LLM.

\section{Hướng phát triển tiếp theo}
\begin{enumerate}[leftmargin=1.6em]
  \item Hoàn thiện backend route và service layer còn thiếu.
  \item Thêm entity resolution tốt hơn cho Author/Organization/Conference.
  \item Xây dựng validation loop cho agent để tăng độ chính xác trả lời.
  \item Thêm bộ test tích hợp cho ingest và query.
  \item Hoàn thiện UI graph với filter theo loại node, chọn subgraph theo paper, và highlight đường đi quan hệ.
  \item Bổ sung quan sát hệ thống: log, metrics, và status pipeline rõ ràng hơn.
\end{enumerate}

\section{Kết luận}
Dự án GraphRAG là một hệ thống có tính thực tiễn cao cho bài toán hỏi đáp trên tài liệu khoa học. Điểm nổi bật của dự án là không chỉ tìm kiếm văn bản, mà còn mô hình hóa tài liệu thành tri thức có cấu trúc để phục vụ suy luận đa bước.

Về mặt kỹ thuật, hệ thống thể hiện sự kết hợp giữa:
\begin{itemize}[leftmargin=1.6em]
  \item xử lý tài liệu khoa học,
  \item LLM extraction,
  \item graph database,
  \item vector search,
  \item và agent reasoning.
\end{itemize}

Nếu tiếp tục hoàn thiện các module còn stub và ổn định hóa extraction, đây có thể là một nền tảng tốt cho nghiên cứu hoặc demo sản phẩm về GraphRAG trong môi trường học thuật.

\end{document}